\documentclass[11pt]{article}

\usepackage[margin=1in]{geometry}
\usepackage{amsmath,amssymb}
\usepackage{mathtools}
\usepackage{microtype}

\title{Parameter Conditions for Improving a Square-Root Bound}
\author{}
\date{}

\begin{document}

\maketitle

\section{Setup}

Let the old function be
\[
H(a)
=
1+\frac{c}{da+\sqrt{a(ea+f)}},
\]
and let the new function be
\[
\widetilde H(a)
=
1+\frac{\widetilde c}
{\widetilde d\,a+\sqrt{a(\widetilde e\,a+\widetilde f)}}.
\]
We consider the interval
\[
0<a\le \frac12.
\]

Throughout, the functions are assumed to be real-valued and to have positive denominators on the relevant interval. Thus the square-root arguments must be nonnegative and the denominators must be strictly positive wherever the inequalities are imposed.

The goal is to choose new parameters
\[
\widetilde c,\ \widetilde d,\ \widetilde e,\ \widetilde f
\]
so that
\[
\boxed{
H(a)<\widetilde H(a)<1+\frac{1}{6a}
\qquad
\text{for all }0<a\le \frac12.
}
\]
Thus the new function should lie strictly above the old one, but strictly below the ceiling
\[
1+\frac{1}{6a}.
\]

\section{Denominator notation}

Define
\[
S(a)
=
da+\sqrt{a(ea+f)}
\]
and
\[
\widetilde S(a)
=
\widetilde d\,a+\sqrt{a(\widetilde e\,a+\widetilde f)}.
\]
Then
\[
H(a)=1+\frac{c}{S(a)}
\]
and
\[
\widetilde H(a)=1+\frac{\widetilde c}{\widetilde S(a)}.
\]

Assuming all denominators are positive, the desired condition is equivalent to
\[
\frac{c}{S(a)}
<
\frac{\widetilde c}{\widetilde S(a)}
<
\frac{1}{6a}
\qquad
\text{for all }0<a\le\frac12.
\]

\section{Exact improvement condition}

The inequality
\[
H(a)<\widetilde H(a)
\]
is equivalent to
\[
\frac{c}{S(a)}
<
\frac{\widetilde c}{\widetilde S(a)}.
\]
Assuming \(c>0\), \(\widetilde c>0\), \(S(a)>0\), and \(\widetilde S(a)>0\), this becomes
\[
\boxed{
\widetilde S(a)
<
\frac{\widetilde c}{c}S(a)
\qquad
\text{for all }0<a\le\frac12.
}
\]
Written out explicitly,
\[
\boxed{
\widetilde d\,a+\sqrt{a(\widetilde e\,a+\widetilde f)}
<
\frac{\widetilde c}{c}
\left(
da+\sqrt{a(ea+f)}
\right)
\qquad
\text{for all }0<a\le\frac12.
}
\]

\section{Exact upper-bound condition}

The condition
\[
\widetilde H(a)<1+\frac{1}{6a}
\]
is equivalent to
\[
\frac{\widetilde c}{\widetilde S(a)}
<
\frac{1}{6a}.
\]
Assuming \(\widetilde c>0\) and \(\widetilde S(a)>0\), this becomes
\[
\boxed{
\widetilde S(a)>6a\widetilde c
\qquad
\text{for all }0<a\le\frac12.
}
\]
Explicitly,
\[
\boxed{
\widetilde d\,a+\sqrt{a(\widetilde e\,a+\widetilde f)}
>
6a\widetilde c
\qquad
\text{for all }0<a\le\frac12.
}
\]

\section{Parameter form of the upper-bound condition}

Set
\[
\widetilde k=6\widetilde c-\widetilde d.
\]
Then the upper-bound condition becomes
\[
\sqrt{a(\widetilde e\,a+\widetilde f)}
>
a\widetilde k.
\]

If
\[
\widetilde k\le 0,
\]
then the right-hand side is nonpositive, while the square-root term is nonnegative.
Thus the upper-bound condition is automatically satisfied whenever the square root is real and the denominator is positive.

If
\[
\widetilde k>0,
\]
then squaring gives the equivalent condition
\[
a(\widetilde e\,a+\widetilde f)
>
a^2\widetilde k^2.
\]
Since \(a>0\), this is equivalent to
\[
\boxed{
\widetilde f+a\left(\widetilde e-\widetilde k^2\right)>0
\qquad
\text{for all }0<a\le\frac12.
}
\]
That is,
\[
\boxed{
\widetilde f+a\left(\widetilde e-(6\widetilde c-\widetilde d)^2\right)>0
\qquad
\text{for all }0<a\le\frac12.
}
\]

\section{Standing reality and positivity conditions}

The comparison inequalities are meaningful only under the following standing assumptions. To keep the functions real-valued on \(0<a\le\frac12\), one needs
\[
\boxed{
ea+f\ge 0
\qquad
\text{and}
\qquad
\widetilde e\,a+\widetilde f\ge 0
\qquad
\text{for all }0<a\le\frac12.
}
\]
For the new function, this means
\[
\boxed{
\widetilde e\,a+\widetilde f\ge 0
\qquad
\text{for all }0<a\le\frac12.
}
\]
The new denominator must also remain positive:
\[
\boxed{
\widetilde d\,a+\sqrt{a(\widetilde e\,a+\widetilde f)}>0
\qquad
\text{for all }0<a\le\frac12.
}
\]

\section{Consequences of assuming \(c<d\)}

Now add the assumption
\[
\boxed{
c<d.
}
\]
For the old function
\[
H(a)
=
1+\frac{c}{da+\sqrt{a(ea+f)}},
\]
this assumption says that the coefficient in the numerator is smaller than the linear coefficient in the denominator.

The ceiling comparison
\[
H(a)<1+\frac{1}{6a}
\]
is equivalent to
\[
\frac{c}{da+\sqrt{a(ea+f)}}<\frac{1}{6a}.
\]
Assuming positive denominators, this is equivalent to
\[
6ac<da+\sqrt{a(ea+f)}.
\]
Thus
\[
\boxed{
\sqrt{a(ea+f)}>a(6c-d).
}
\]
Set
\[
k=6c-d.
\]
Then the ceiling condition becomes
\[
\boxed{
\sqrt{a(ea+f)}>ak.
}
\]

\subsection*{Case 1: \(d\ge 6c\)}

If
\[
\boxed{
d\ge 6c,
}
\]
then
\[
k=6c-d\le 0.
\]
The right-hand side \(ak\) is nonpositive for \(a>0\), while the square-root term is nonnegative.
Therefore the ceiling condition is automatic, up to the already assumed real-valuedness and positivity conditions.

More precisely, if \(d>6c\), then
\[
da>6ac,
\]
so
\[
da+\sqrt{a(ea+f)}>6ac
\]
for every \(a>0\).
Hence
\[
\boxed{
H(a)<1+\frac{1}{6a}
\qquad
\text{for all }0<a\le\frac12.
}
\]

If \(d=6c\), then the inequality becomes
\[
\sqrt{a(ea+f)}>0.
\]
Thus strict inequality holds whenever
\[
ea+f>0.
\]
If the square-root term vanishes at some point, then the old function touches the ceiling at that point.

\subsection*{Case 2: \(c<d<6c\)}

If
\[
\boxed{
c<d<6c,
}
\]
then
\[
k=6c-d>0.
\]
In this case, \(c<d\) alone is not enough to guarantee that \(H\) lies below the ceiling.
We must impose
\[
\sqrt{a(ea+f)}>a(6c-d).
\]
Squaring is valid because both sides are nonnegative.
For \(a>0\), this gives
\[
a(ea+f)>a^2(6c-d)^2,
\]
or equivalently
\[
\boxed{
f+a\left(e-(6c-d)^2\right)>0
\qquad
\text{for all }0<a\le\frac12.
}
\]

Therefore, under the assumption \(c<d<6c\), the old function remains strictly below the ceiling exactly when
\[
\boxed{
f+a\left(e-(6c-d)^2\right)>0
}
\]
throughout the interval, together with the standing real-valuedness and positivity assumptions.

\subsection*{Endpoint contact}

If
\[
f+a\left(e-(6c-d)^2\right)=0
\]
at some point \(a\), then
\[
H(a)=1+\frac{1}{6a}.
\]
Thus equality with the ceiling occurs exactly at the points where
\[
\boxed{
a=
\frac{f}{(6c-d)^2-e},
}
\]
provided \(6c-d>0\) and the displayed value lies in the interval under consideration.

\subsection*{Application to new parameters}

The same analysis applies to the new function after replacing
\[
c,d,e,f
\]
by
\[
\widetilde c,\widetilde d,\widetilde e,\widetilde f.
\]
Thus, if
\[
\boxed{
\widetilde d\ge 6\widetilde c,
}
\]
then the new function automatically stays below the ceiling, subject to real-valuedness and positivity.
If instead
\[
\boxed{
\widetilde c<\widetilde d<6\widetilde c,
}
\]
then the required ceiling condition is
\[
\boxed{
\widetilde f
+
a\left(\widetilde e-(6\widetilde c-\widetilde d)^2\right)>0
\qquad
\text{for all }0<a\le\frac12.
}
\]

\section{A clean sufficient strategy}

A simple sufficient way to make the new function larger than the old one is to increase the numerator and reduce the denominator.
For example, suppose
\[
\boxed{
\widetilde c\ge c,
}
\]
and
\[
\boxed{
\widetilde d\le d,\qquad
\widetilde e\le e,\qquad
\widetilde f\le f,
}
\]
with at least one inequality strict in a way that makes
\[
\frac{\widetilde c}{\widetilde S(a)}
>
\frac{c}{S(a)}
\]
on the full interval.

This is only a sufficient strategy, not a necessary one.
After making such a choice, one must still impose the ceiling condition
\[
\boxed{
\widetilde d\,a+\sqrt{a(\widetilde e\,a+\widetilde f)}
>
6a\widetilde c
\qquad
\text{for all }0<a\le\frac12.
}
\]
This prevents the new function from crossing above
\[
1+\frac{1}{6a}.
\]

\section{Specific initial choice: \(c=2,\ d=3,\ e=49,\ f=16\)}

Now suppose the initial parameter choice is
\[
\boxed{
c=2,\qquad d=3,\qquad e=49,\qquad f=16.
}
\]
Then the old function is
\[
H(a)
=
1+\frac{2}{3a+\sqrt{a(49a+16)}}.
\]
Thus
\[
S(a)
=
3a+\sqrt{a(49a+16)}.
\]

The exact improvement and ceiling conditions for the new parameters become
\[
\boxed{
6a\widetilde c
<
\widetilde d\,a+\sqrt{a(\widetilde e\,a+\widetilde f)}
<
\frac{\widetilde c}{2}
\left(
3a+\sqrt{a(49a+16)}
\right)
}
\]
for every \(a\) in the interval under consideration.

\subsection*{Endpoint behavior}

Before choosing new parameters, one should check whether the old function is already strictly below the ceiling
\[
1+\frac{1}{6a}.
\]
The condition
\[
H(a)<1+\frac{1}{6a}
\]
is equivalent to
\[
\frac{2}{3a+\sqrt{a(49a+16)}}<\frac{1}{6a}.
\]
Therefore
\[
3a+\sqrt{a(49a+16)}>12a,
\]
or
\[
\sqrt{a(49a+16)}>9a.
\]
Squaring, for \(a>0\), gives
\[
49a^2+16a>81a^2,
\]
which is equivalent to
\[
16a>32a^2.
\]
Hence
\[
\boxed{
H(a)<1+\frac{1}{6a}
\quad\Longleftrightarrow\quad
0<a<\frac12.
}
\]
At the endpoint \(a=\frac12\), the old function touches the ceiling:
\[
\boxed{
H\left(\frac12\right)
=
1+\frac{1}{6\cdot\frac12}.
}
\]

If touching the ceiling at the endpoint is allowed, then this is acceptable boundary behavior.
The strict comparison should then be imposed on the open interval
\[
\boxed{
0<a<\frac12.
}
\]
At the endpoint \(a=\frac12\), one may separately impose the non-strict ceiling condition
\[
\boxed{
\widetilde H\left(\frac12\right)
\le
1+\frac{1}{6\cdot\frac12}.
}
\]
In denominator form, this endpoint condition is
\[
\boxed{
\widetilde S\left(\frac12\right)\ge 3\widetilde c.
}
\]
Equivalently,
\[
\boxed{
\frac{\widetilde d}{2}
+
\sqrt{\frac12\left(\frac{\widetilde e}{2}+\widetilde f\right)}
\ge
3\widetilde c.
}
\]

If one instead requires strict improvement
\(\widetilde H(a)>H(a)\) also at \(a=\frac12\), while also requiring
\(\widetilde H(a)\le 1+\frac{1}{6a}\) at \(a=\frac12\), then no such new function can exist, because \(H(\frac12)\) is already equal to the ceiling.

\subsection*{Conditions for the new parameters on \(0<a<\frac12\)}

On the open interval \(0<a<\frac12\), the new parameters must satisfy
\[
\boxed{
6a\widetilde c
<
\widetilde d\,a+\sqrt{a(\widetilde e\,a+\widetilde f)}
<
\frac{\widetilde c}{2}
\left(
3a+\sqrt{a(49a+16)}
\right)
}
\]
for every \(0<a<\frac12\).

Equivalently, writing
\[
\widetilde S(a)
=
\widetilde d\,a+\sqrt{a(\widetilde e\,a+\widetilde f)},
\]
the condition is
\[
\boxed{
6a\widetilde c
<
\widetilde S(a)
<
\frac{\widetilde c}{2}S(a)
\qquad
\text{for all }0<a<\frac12.
}
\]

The left inequality is the ceiling condition. With
\[
\widetilde k=6\widetilde c-\widetilde d,
\]
it can be written as
\[
\sqrt{a(\widetilde e\,a+\widetilde f)}
>
a\widetilde k.
\]
If
\[
\widetilde k\le 0,
\]
then this condition is automatic whenever the square root is real and the denominator is positive.
If
\[
\widetilde k>0,
\]
then it is equivalent to
\[
\boxed{
\widetilde f
+
a\left(\widetilde e-(6\widetilde c-\widetilde d)^2\right)
>0
\qquad
\text{for all }0<a<\frac12.
}
\]

The right inequality is the improvement condition:
\[
\boxed{
\widetilde d\,a+\sqrt{a(\widetilde e\,a+\widetilde f)}
<
\frac{\widetilde c}{2}
\left(
3a+\sqrt{a(49a+16)}
\right)
\qquad
\text{for all }0<a<\frac12.
}
\]
This says that the new denominator must be small enough, relative to \(\widetilde c\), to make the new function larger than the old one.

\subsection*{Real-valuedness}

The new square root must be real on the interval, so one also needs
\[
\boxed{
\widetilde e\,a+\widetilde f\ge 0
\qquad
\text{for all }0<a<\frac12.
}
\]
The new denominator should remain positive:
\[
\boxed{
\widetilde d\,a+\sqrt{a(\widetilde e\,a+\widetilde f)}>0
\qquad
\text{for all }0<a<\frac12.
}
\]

\section{Summary}

The exact conditions for
\[
H(a)<\widetilde H(a)<1+\frac{1}{6a}
\]
on the interval \(0<a\le\frac12\) are:
\[
\boxed{
\widetilde d\,a+\sqrt{a(\widetilde e\,a+\widetilde f)}
<
\frac{\widetilde c}{c}
\left(
da+\sqrt{a(ea+f)}
\right)
}
\]
and
\[
\boxed{
\widetilde d\,a+\sqrt{a(\widetilde e\,a+\widetilde f)}
>
6a\widetilde c
}
\]
for every \(0<a\le\frac12\), together with the necessary real-valuedness and positivity assumptions.

Equivalently, the new denominator must satisfy
\[
\boxed{
6a\widetilde c
<
\widetilde S(a)
<
\frac{\widetilde c}{c}S(a)
\qquad
\text{for all }0<a\le\frac12.
}
\]

\end{document}
